\documentclass[10pt,a4paper]{amsart}
\usepackage[margin=19mm]{geometry}
\usepackage{amsmath,amssymb,mathtools}
\usepackage[scr=rsfs,cal=euler]{mathalfa}
\usepackage{xfrac}
\usepackage[unicode,hidelinks,bookmarks=false]{hyperref}
\newcommand{\ie}{i.e.\ }
\newcommand{\cf}{cf.\ }
\newcommand{\resp}{respectively\ }
\newcommand{\Subsection}{\subsection}
\providecommand{\nobreakdash}{\nobreak\hbox{-}}
\setlength{\emergencystretch}{2em}
\multlinegap0pt
\usepackage{amssymb}
\usepackage{xcolor}
\newtheorem{lemma}{Lemma}
\title{Critical window for approximate counting\\ in dense Ising models}
\author{Andreas Galanis, Daniel Stefankovic, Eric Vigoda}
\date{Source: arXiv:2603.21406v1 (CC BY 4.0)}
\begin{document}
\maketitle

\noindent\emph{Source and licence.} Critical window for approximate counting\\ in dense Ising models. Version arXiv:2603.21406v1 (CC BY 4.0), distributed by arXiv under the Creative Commons Attribution 4.0 International licence, http://creativecommons.org/licenses/by/4.0/, as recorded in the arXiv licence metadata for this exact version and shown on the abstract page https://arxiv.org/abs/2603.21406v1. Full licence notice: \texttt{licenses/arxiv-2603.21406-CC-BY-4.0.txt}.\par\smallskip

\noindent\textit{Editorial scope.} Counted unit: the article's lemma lem:maxb (source0.tex lines 214--216). The article's complete original proof of it is delivered with it (source0.tex lines 217--243, one \texttt{proof} environment of 27 lines). No mathematics is written, repaired or re-derived: the statement and the proof are the article's own lines, in the article's own notation, and no retained line is substituted or re-flowed.  Retained uncounted as the source-local context the proof needs: lines 183--185; lines 196--198.  Nothing was omitted from any retained span, so the package carries no omission record.  No retained line contains a citation, so the wrapper carries no bibliography and no bibliography entry.  No retained line contains a figure environment, an image-inclusion command or drawing code, so no image is delivered and the package declares no asset dependency.  Editorial formatting only, in addition: an AMS \texttt{amsart} preamble that restates the article's own declarations and macros used by the retained lines, namely the environments \texttt{lemma} on separate counters, together with the standard packages those lines need.  The retained environments print in this document as Lemma 1 in order of appearance, whereas the article prints the same items with the numbering of its own full document.  The complete native source of the article as distributed in the arXiv e-print for this exact version is bundled unchanged as \texttt{sources/196943/source0.tex}.
\begin{equation}\label{ZOB}
Q(b):= 2\beta b^2 + t\, H\big(\tfrac{1}{2} + \tfrac{b}{t}\big) \geq  2\beta b^2 + \log {t\choose t/2 + b},
\end{equation}

\begin{equation}\label{e1}
\beta = \frac{1}{4\hat{b}}\ln \frac{t + 2\hat{b}}{t - 2\hat{b}}.
\end{equation}

\begin{lemma}\label{lem:maxb}
For any $\hat{b}\in (0,t)$ and $\beta=\frac{1}{4\hat{b}}\ln \frac{t + 2\hat{b}}{t - 2\hat{b}}$ as in~\eqref{e1}, the maximum of $Q(b)$ in \eqref{ZOB} over $b\in [0,t/2]$  is achieved uniquely at $b=\hat{b}$.
\end{lemma}

\begin{proof}
We have
\begin{equation}\label{epart}
\frac{\partial}{\partial b} Q(b) = 4\beta b - \ln\frac{t+2b}{t-2b}.
\end{equation}
For~\eqref{epart} to be zero we need to have
$$\beta t = \frac{1}{2} \frac{t}{2b} \ln\frac{1+2b/t}{1-2b/t}.$$
We will show that the function
\begin{equation}\label{eeo}
y\mapsto \frac{1}{y}\ln\frac{1+y}{1-y}
\end{equation}
is increasing on $[0,1]$ and hence there is unique critical $b$ (and we chose $\beta$ so that $\hat{b}$ is critical).

The derivative of~\eqref{eeo} is
$$
\frac{2}{y(1-y^2)} - \frac{1}{y^2}\ln\frac{1+y}{1-y},
$$
which we will show is positive on $(0,1)$. It is sufficient to show that the following is increasing
\begin{equation}\label{eeozz}
\frac{2y}{(1-y^2)} - \ln\frac{1+y}{1-y},
\end{equation}
since the value at $y=0$ is $0$. The derivative of~\eqref{eeozz} 
$$
\frac{4y^2}{(1-y^2)^2}>0,
$$
true for $y\in (0,1)$.
\end{proof}

\end{document}
